\section{Control simultáneo de todas las subactivaciones}
Sea $z_e(Y)={\bf1}_{e\in Y}$, $z_S=\prod_{e\in S}z_e$. Toda función del cubo finito
se escribe $f=\sum_S f_S z_S$ con $f_S=[S]f$. Dotamos a este álgebra de la norma
$\|f\|_\rho=\sum_S |f_S|\rho^{|S|}$. Como $z_S z_T=z_{S\cup T}$ y $\rho\ge1$,
la norma es submultiplicativa. Cada evaluación real tiene norma a lo sumo uno.
Para vectores usamos máximo de normas y para matrices la máxima suma por fila de
normas de entradas. Esta elección evita un coste por el número de subactivaciones.

Sea $\Pi(H)$ el logaritmo por entrada del escalado de $P\exp H$ a márgenes uno
respecto de $P$, en el soporte. Lo construimos en el álgebra: se buscan potenciales
$a,b$ en el gauge $\sum a_i=\sum b_i$, tales que $P_{ij}\exp(H_{ij}+a_i+b_j)$
tenga márgenes uno y se pone $\Pi(H)_{ij}=H_{ij}+a_i+b_j$. El defecto de márgenes
vive en el complemento del gauge nulo $({\bf1},-{\bf1})$. En cero el jacobiano es
$D_0=\left(\begin{smallmatrix}I&P\\P^t&I\end{smallmatrix}\right)$.
Su inverso en ese complemento tiene norma uniforme de sumas absolutas: la brecha da
$\|D_0^+ y\|_2\le C\sqrt{2j}\|y\|_\infty$, y en la ecuación de una coordenada la
fila cruzada tiene norma $\ell^2$ a lo sumo $\sqrt{M_0/j}$. Así
$\|D_0^+y\|_\infty\le C'\|y\|_\infty$ para datos compatibles.
Proyectar un dato cualquiera sobre el complemento cuesta a lo sumo dos en supremo.
La cota escalar de sumas por fila se transfiere al álgebra por desigualdad triangular.
Las ecuaciones tienen resto de Taylor cuadrático, con constante uniforme: cada
margen es suma con pesos positivos de suma uno y la exponencial tiene mayorante
absoluta en una bola fija. La iteración con el inverso basal tiene constante menor
que $1/2$ reduciendo la bola. Restar dos ecuaciones muestra que $\Pi(0)=0$ y que
$\Pi$ es Lipschitz con constante uniforme en un entorno fijo de $H=0$.

Pongamos $q_e=P_e\exp v_e$, $h_f=1-\sum_{e\in f}z_e q_e$ y $d_e=2+2z_e$.
La ecuación de estacionariedad se expresa
\[
 v=\Pi(H(v)),\qquad
 H_e=z_e\sum_{f\ni e}\log h_f-(1+2z_e)\log(1-q_e)+\log(1-P_e),
\]
donde la suma tiene las dos pendientes. Se obtiene al dividir la ecuación para
$q_e(1-q_e)^{d_e-1}$ por la de torres para $P_e(1-P_e)$, con los potenciales RC.
Al evaluar en una máscara, las fibras unitarias son equivalentes a no imponerlas.
En una bola $\|v\|_\infty\le a$ de radio fijo pequeño,
$\|1-h_f\|_\rho\le\rho e^a\epsilon_A$ y $\|q_e\|_\rho\le e^a P_e$.
Las series de logaritmos convergen. La norma de $H(0)$ y la constante de Lipschitz
de $H$ son $O_\rho(\epsilon_A+M_0/j)$: el diferencial de una carga suma como máximo
$\rho e^a\epsilon_A\|\Delta v\|_\infty$, y el de un término local cuesta $O_\rho(P_e)$.
Con $\epsilon_0$ pequeño y $j_0$ grande, el mapa compuesto es contractivo en una
bola invariante de radio $C(\epsilon_A+1/j)$. Se obtiene
$\max_e\|v_e\|_\rho\le C(\epsilon_A+1/j)$, sin dependencia en $|Y|$.

En cada evaluación real, $q$ es positivo, satisface márgenes, $h_f>0$ y $\max q_e<e^{-4}$
aumentando $j_0$. Es el máximo global único: en el sitio de grado $d\in\{2,4\}$,
$s_d(x)=-x\log x+(d-1)(1-x)\log(1-x)$ es cóncava hasta $1/d$.
La tangente en $q<e^{-4}$ domina estrictamente salvo en $q$; desde $1/d$ hasta uno
la diferencia tangente menos función es cóncava y en uno vale
$-\log q-d(1-q)>0$. Se añaden los términos cóncavos $-h_f\log h_f$ y se usa
la estacionariedad en la clase factible. Esta identificación emplea el soporte y
los mismos márgenes de torres en todas las evaluaciones.
